\section{Metodología Experimental}

\subsection{Procedimiento seguido para la toma de datos}

La práctica experimental se desarrolló en cuatro etapas, cada una orientada a una función específica del osciloscopio.

\begin{enumerate}
    \item \textbf{Osciloscopio como voltímetro de corriente continua (CC).} Se conectó el barrido horizontal y se fijó el selector AC-GND-DC en la posición DC. A continuación, se centró la traza sobre la pantalla del osciloscopio y se aplicó al canal 1 una tensión continua variable. Para cada valor, se registró la lectura del multímetro digital y se contó el número de subdivisiones verticales recorridas por la traza, con el fin de calcular el valor experimental del voltaje mediante la expresión
    \[
    V_{DC}^{exp} = \frac{nG}{5},
    \]
    donde $n$ es el número de divisiones verticales y $G$ la ganancia vertical seleccionada en V/div. Este procedimiento se repitió para cuatro valores distintos de tensión.

    \item \textbf{Osciloscopio como voltímetro de corriente alterna (CA).} Se cambió el selector AC-GND-DC a la posición AC y se aplicó una señal alterna al canal 1. La tensión se midió en paralelo con un voltímetro de CA, cuya lectura eficaz se tomó como referencia. Se ajustaron el tiempo de barrido y el control de sincronismo para obtener una forma de onda estable en pantalla. A partir del número de subdivisiones verticales, se determinó la amplitud pico a pico $V_{PP}$ y se calculó el valor teórico del voltaje eficaz como
    \[
    V_{rms}^{teo} = \frac{V_{PP}}{2\sqrt{2}}.
    \]
    Este proceso se repitió para cinco valores de tensión alterna distintos.

    \item \textbf{Osciloscopio como frecuencímetro.} Se aplicó al canal 1 una señal sinusoidal procedente del generador de frecuencias y se ajustó el barrido hasta que una oscilación completa quedara visible en la pantalla. Para cada valor de frecuencia, se midió el periodo experimental a partir del número de divisiones horizontales ocupadas por un ciclo completo:
    \[
    T_{exp} = \frac{nE}{5},
    \]
    donde $n$ es el número de divisiones horizontales y $E$ la escala de tiempo en ms/div. La frecuencia experimental se obtuvo como la inversa del periodo y se comparó con la frecuencia teórica indicada por el generador. Se repitió la medición para cinco frecuencias distintas comprendidas entre 60 Hz y 1000 Hz.

    \item \textbf{Osciloscopio como medidor de desfase.} Se montó un circuito RC con $R = 1{,}99\,\mathrm{k\Omega}$ y $C = 95{,}1\,\mathrm{nF}$, alimentado con una señal senoidal de 10 Vpp. La señal de entrada se conectó al canal 1 y la señal del capacitor al canal 2. Se ajustó el osciloscopio para visualizar dos ondas completas en pantalla y se midieron las subdivisiones horizontales correspondientes al desfase, $m$, y al período completo, $n$. El ángulo de desfase experimental se calculó mediante
    \[
    \phi_{exp} = 360^\circ \left(\frac{m}{n}\right).
    \]
    La frecuencia de la señal se varió entre 300 Hz y 600 Hz en incrementos de 25 Hz. Para cada valor, se obtuvo el desfase y, posteriormente, el valor experimental de la constante de tiempo $\tau$ mediante una regresión lineal de $\tan(\phi)$ en función del periodo $T$, utilizando la relación
    \[
    \tan(\phi) = \omega \tau = \frac{2\pi\tau}{T}.
    \]

\end{enumerate}

\subsection{Variables medidas directamente}

Las variables medidas directamente durante la experiencia fueron:

\begin{itemize}
    \item Tensión de corriente continua y alterna con el multímetro digital.
    \item Número de subdivisiones verticales y horizontales observadas en la pantalla del osciloscopio.
    \item Ganancia vertical $G$ (V/div) y escala de tiempo $E$ (ms/div) seleccionadas.
\end{itemize}

\subsection{Variables obtenidas indirectamente}

A partir de las mediciones gráficas y de la relación entre la escala del osciloscopio y la señal estudiada, se obtuvieron las siguientes magnitudes:

\begin{itemize}
    \item $V_{DC}^{exp}$ y $V_{PP}^{exp}$, a partir de las subdivisiones verticales y la ganancia.
    \item $V_{rms}^{teo}$, a partir de $V_{PP}$.
    \item El período y la frecuencia experimental, a partir de la escala de tiempo y el número de divisiones horizontales.
    \item El ángulo de desfase $\phi_{exp}$ y la constante de tiempo $\tau_{exp}$, mediante la regresión lineal de $\tan(\phi)$ frente a $T$.
\end{itemize}

\subsection{Datos asumidos conocidos}

Se consideraron conocidos los siguientes parámetros:

\begin{itemize}
    \item Los valores nominales de los componentes del circuito RC: $R = 1{,}99\,\mathrm{k\Omega}$ y $C = 95{,}1\,\mathrm{nF}$.
    \item La frecuencia indicada por el generador de señales, que se usó como referencia teórica.
    \item Los rangos de precisión y resolución de los instrumentos utilizados.
\end{itemize}

\subsection{Limitaciones del diseño experimental}

El procedimiento experimental permite comparar directamente la lectura gráfica del osciloscopio con instrumentos de referencia, como el multímetro digital y el generador de señales. Esta estrategia facilita la evaluación de la precisión del osciloscopio en distintos regímenes de operación. Sin embargo, la lectura de las subdivisiones en pantalla depende fuertemente de la apreciación visual del operador, lo que introduce subjetividad en las mediciones. Esto resulta especialmente relevante en la determinación del desfase, donde pequeños errores en la lectura de $m$ y $n$ afectan significativamente el valor calculado del ángulo.

\subsection{Fuentes de error más críticas}

Las principales fuentes de error del experimento fueron:

\begin{itemize}
    \item La resolución visual para contar subdivisiones en la pantalla del osciloscopio, incluida la dificultad para estimar fracciones de división.
    \item La estabilidad de la sincronización de la señal, que puede impedir fijar una forma de onda completamente estática.
    \item La tolerancia de los componentes $R$ y $C$, que afecta al valor teórico de $\tau$ calculado.
\end{itemize}
